% =============================================================================
% embedded_mindset.tex – Fragment für techdoc.cls
% =============================================================================
% Einbinden in main.tex mit: % =============================================================================
% embedded_mindset.tex – Fragment für techdoc.cls
% =============================================================================
% Einbinden in main.tex mit: % =============================================================================
% embedded_mindset.tex – Fragment für techdoc.cls
% =============================================================================
% Einbinden in main.tex mit: % =============================================================================
% embedded_mindset.tex – Fragment für techdoc.cls
% =============================================================================
% Einbinden in main.tex mit: \input{content/embedded_mindset}
% =============================================================================

\section{Desktop → Embedded Mindset (10 Regeln)}
\label{sec:embedded-mindset}

\begin{enumerate}
  \item \textbf{Hardware ist die API}\\
    Pins, Busse, Register sind \glqq die Schnittstelle\grqq. OS-Abstraktionen fehlen oft.

  \item \textbf{Worst-Case zählt, nicht der Durchschnitt}\\
    Plane Deadlines/Latenzen, nicht \glqq läuft meistens\grqq.

  \item \textbf{Kein Blockieren im kritischen Pfad}\\
    Keine langen ISR, keine Busy-Waits, kein \code{delay()} in Logik/Tasks.

  \item \textbf{ISR = minimal}\\
    ISR setzt Flag/Queue + Timestamp. Verarbeitung im Task/Mainloop.

  \item \textbf{Speicher ist ein Budget}\\
    Stack/RAM sind eng. Große Daten statisch oder in Pools. Rekursion vermeiden.

  \item \textbf{Nebenläufigkeit ist überall}\\
    ISR ↔ Task ↔ Main: Ownership und Synchronisation (Queue/Atomics) festlegen.

  \item \textbf{C++ Feature-Set bewusst wählen}\\
    \code{constexpr}, \code{enum class}, kleine POD-Structs: ja. Exceptions/RTTI/iostream: meist nein.

  \item \textbf{I/O entkoppeln}\\
    \code{read\_inputs()} → \code{logic\_update()} → \code{write\_outputs()} (testbar, deterministisch).

  \item \textbf{Fehler sind Normalfall}\\
    Timeouts, Brownout, Disconnect: Recovery + Watchdog einplanen.

  \item \textbf{Observability ist knapp}\\
    Logging sparsam, Counters, Debug-Pins/Tracing. Debugbarkeit mitdesignen.
\end{enumerate}

% -----------------------------------------------------------------------------
\subsection{Beispiel: Selection Panel (One-Hot) -- \glqq Taster i wählt LED i\grqq}
\label{sec:selection-panel-beispiel}

\textbf{Ziel:} 10~Taster, 10~LEDs. Bei Tastendruck wird \highlight{genau eine} LED aktiv (One-Hot). Ein neuer Tastendruck schaltet \textbf{sofort} auf die neue Auswahl um.

\subsubsection{Konventionen (wichtig, konsistent)}

\begin{itemize}
  \item \textbf{Buttons sind active-low (Hardware):} \code{raw\_bit = 0} bedeutet \glqq gedrückt\grqq{} (\zB CD4021).
  \item \textbf{Logik arbeitet normalisiert:} \code{pressed = !raw\_bit} → dann gilt in der Logik: \code{1 = gedrückt}.
  \item \textbf{Auswahl wird als ID geführt:}\\
    \code{selected\_id} $\in \{0..10\}$\\
    \code{0} = keine Auswahl, \code{1..10} = LED/Taster-ID
\end{itemize}

\subsubsection{Datenfluss (Architekturregel)}

\begin{lstlisting}[style=base,numbers=none]
read_inputs_raw() -> normalize_active_low() -> debounce_update()
                  -> selection_update() -> write_outputs_onehot()
\end{lstlisting}

\begin{table}[H]
\centering
\caption{Datenfluss-Funktionen im Detail}
\label{tab:datenfluss-funktionen}
\begin{tabularx}{\textwidth}{l X}
\toprule
\textbf{Funktion} & \textbf{Beschreibung} \\
\midrule
\code{read\_inputs\_raw()}      & Liest Bits aus dem Schieberegister (\zB 2~Bytes für 10~Taster) \\
\code{normalize\_active\_low()} & Invertiert Bits → Logik-Semantik \code{1 = pressed} \\
\code{debounce\_update()}       & Macht aus prellenden Rohwerten stabile Zustände \\
\code{selection\_update()}      & Setzt \code{selected\_id} bei neuem Press-Event \\
\code{write\_outputs\_onehot()} & Setzt genau \textbf{ein} LED-Bit (alle anderen 0) \\
\bottomrule
\end{tabularx}
\end{table}

% -----------------------------------------------------------------------------
\subsection{Arduino-Loop vs.\ FreeRTOS (Prinzip, nicht Code)}
\label{sec:loop-vs-rtos}

\subsubsection{Variante A: Arduino-Loop (ohne eigene Tasks)}

\code{loop()} läuft zyklisch (Polling), \zB alle \SI{1}{\milli\second}. Alle Schritte passieren sequentiell:

\begin{lstlisting}[style=base,numbers=none]
Inputs lesen -> normalisieren -> entprellen -> Event erkennen -> LED schreiben
\end{lstlisting}

\begin{table}[H]
\centering
\caption{Arduino-Loop: Stärken und Risiken}
\label{tab:arduino-loop}
\begin{tabularx}{\textwidth}{l X}
\toprule
\textbf{Stärke} & Simpel, wenige Nebenläufigkeitsprobleme \\
\textbf{Risiko} & \code{loop()} darf nicht blockieren (Regel~3) \\
\bottomrule
\end{tabularx}
\end{table}

\subsubsection{Variante B: FreeRTOS (Task-basiert)}

\begin{itemize}
  \item \textbf{IO-Task (high prio)} besitzt SPI und macht:\\
    Inputs lesen → normalisieren → entprellen → \code{selected\_id} bestimmen
  \item Optional: \textbf{App-Task} bekommt Events per Queue (\zB für Serial/Web)
\end{itemize}

\begin{table}[H]
\centering
\caption{FreeRTOS: Stärken und Risiken}
\label{tab:freertos}
\begin{tabularx}{\textwidth}{l X}
\toprule
\textbf{Stärke} & Klare Prioritäten, gute Trennung IO ↔ App \\
\textbf{Risiko} & Ownership/Synchronisation sauber halten (Regel~6) \\
\bottomrule
\end{tabularx}
\end{table}

% -----------------------------------------------------------------------------
\subsection{One-Hot-Ausgabe (Definition, byte-basiert)}
\label{sec:one-hot-ausgabe}

\textbf{Regel:} Erst löschen, dann genau 1~Bit setzen.

\begin{lstlisting}[style=arduino]
// Pseudo-Code
clear(led_bytes);
if (selected_id != 0) {
    uint8_t bit = selected_id - 1;
    led_bytes[bit / 8] |= (1u << (bit % 8));
}
\end{lstlisting}

\begin{infobox}[Passend zum verifizierten Mapping]
\begin{itemize}
  \item \code{OUT[0]}: LED1..LED8 → Bit0..Bit7
  \item \code{OUT[1]}: LED9..LED10 → Bit0..Bit1
\end{itemize}
\end{infobox}

% -----------------------------------------------------------------------------
\subsection{One-Hot-Checkliste (damit es wirklich immer stimmt)}
\label{sec:one-hot-checkliste}

\begin{itemize}
  \item[$\square$] \code{selected\_id} ist \textbf{ein} Wert (\code{0..10}), kein Bitfeld
  \item[$\square$] Active-low wird \textbf{einmal} normalisiert (\code{pressed = !raw\_bit})
  \item[$\square$] Event ist \textbf{Press-Edge} (released → pressed), nicht \glqq level\grqq
  \item[$\square$] LED-Ausgabe hat \textbf{einen zentralen Pfad}: \code{clear} → \code{set one bit}
  \item[$\square$] Reset/Boot setzt deterministischen Zustand (\zB \code{selected\_id=0}, alle LEDs aus)
  \item[$\square$] Debug/Logs: nur Events loggen (nicht jeden Scan)
\end{itemize}

% =============================================================================

\section{Desktop → Embedded Mindset (10 Regeln)}
\label{sec:embedded-mindset}

\begin{enumerate}
  \item \textbf{Hardware ist die API}\\
    Pins, Busse, Register sind \glqq die Schnittstelle\grqq. OS-Abstraktionen fehlen oft.

  \item \textbf{Worst-Case zählt, nicht der Durchschnitt}\\
    Plane Deadlines/Latenzen, nicht \glqq läuft meistens\grqq.

  \item \textbf{Kein Blockieren im kritischen Pfad}\\
    Keine langen ISR, keine Busy-Waits, kein \code{delay()} in Logik/Tasks.

  \item \textbf{ISR = minimal}\\
    ISR setzt Flag/Queue + Timestamp. Verarbeitung im Task/Mainloop.

  \item \textbf{Speicher ist ein Budget}\\
    Stack/RAM sind eng. Große Daten statisch oder in Pools. Rekursion vermeiden.

  \item \textbf{Nebenläufigkeit ist überall}\\
    ISR ↔ Task ↔ Main: Ownership und Synchronisation (Queue/Atomics) festlegen.

  \item \textbf{C++ Feature-Set bewusst wählen}\\
    \code{constexpr}, \code{enum class}, kleine POD-Structs: ja. Exceptions/RTTI/iostream: meist nein.

  \item \textbf{I/O entkoppeln}\\
    \code{read\_inputs()} → \code{logic\_update()} → \code{write\_outputs()} (testbar, deterministisch).

  \item \textbf{Fehler sind Normalfall}\\
    Timeouts, Brownout, Disconnect: Recovery + Watchdog einplanen.

  \item \textbf{Observability ist knapp}\\
    Logging sparsam, Counters, Debug-Pins/Tracing. Debugbarkeit mitdesignen.
\end{enumerate}

% -----------------------------------------------------------------------------
\subsection{Beispiel: Selection Panel (One-Hot) -- \glqq Taster i wählt LED i\grqq}
\label{sec:selection-panel-beispiel}

\textbf{Ziel:} 10~Taster, 10~LEDs. Bei Tastendruck wird \highlight{genau eine} LED aktiv (One-Hot). Ein neuer Tastendruck schaltet \textbf{sofort} auf die neue Auswahl um.

\subsubsection{Konventionen (wichtig, konsistent)}

\begin{itemize}
  \item \textbf{Buttons sind active-low (Hardware):} \code{raw\_bit = 0} bedeutet \glqq gedrückt\grqq{} (\zB CD4021).
  \item \textbf{Logik arbeitet normalisiert:} \code{pressed = !raw\_bit} → dann gilt in der Logik: \code{1 = gedrückt}.
  \item \textbf{Auswahl wird als ID geführt:}\\
    \code{selected\_id} $\in \{0..10\}$\\
    \code{0} = keine Auswahl, \code{1..10} = LED/Taster-ID
\end{itemize}

\subsubsection{Datenfluss (Architekturregel)}

\begin{lstlisting}[style=base,numbers=none]
read_inputs_raw() -> normalize_active_low() -> debounce_update()
                  -> selection_update() -> write_outputs_onehot()
\end{lstlisting}

\begin{table}[H]
\centering
\caption{Datenfluss-Funktionen im Detail}
\label{tab:datenfluss-funktionen}
\begin{tabularx}{\textwidth}{l X}
\toprule
\textbf{Funktion} & \textbf{Beschreibung} \\
\midrule
\code{read\_inputs\_raw()}      & Liest Bits aus dem Schieberegister (\zB 2~Bytes für 10~Taster) \\
\code{normalize\_active\_low()} & Invertiert Bits → Logik-Semantik \code{1 = pressed} \\
\code{debounce\_update()}       & Macht aus prellenden Rohwerten stabile Zustände \\
\code{selection\_update()}      & Setzt \code{selected\_id} bei neuem Press-Event \\
\code{write\_outputs\_onehot()} & Setzt genau \textbf{ein} LED-Bit (alle anderen 0) \\
\bottomrule
\end{tabularx}
\end{table}

% -----------------------------------------------------------------------------
\subsection{Arduino-Loop vs.\ FreeRTOS (Prinzip, nicht Code)}
\label{sec:loop-vs-rtos}

\subsubsection{Variante A: Arduino-Loop (ohne eigene Tasks)}

\code{loop()} läuft zyklisch (Polling), \zB alle \SI{1}{\milli\second}. Alle Schritte passieren sequentiell:

\begin{lstlisting}[style=base,numbers=none]
Inputs lesen -> normalisieren -> entprellen -> Event erkennen -> LED schreiben
\end{lstlisting}

\begin{table}[H]
\centering
\caption{Arduino-Loop: Stärken und Risiken}
\label{tab:arduino-loop}
\begin{tabularx}{\textwidth}{l X}
\toprule
\textbf{Stärke} & Simpel, wenige Nebenläufigkeitsprobleme \\
\textbf{Risiko} & \code{loop()} darf nicht blockieren (Regel~3) \\
\bottomrule
\end{tabularx}
\end{table}

\subsubsection{Variante B: FreeRTOS (Task-basiert)}

\begin{itemize}
  \item \textbf{IO-Task (high prio)} besitzt SPI und macht:\\
    Inputs lesen → normalisieren → entprellen → \code{selected\_id} bestimmen
  \item Optional: \textbf{App-Task} bekommt Events per Queue (\zB für Serial/Web)
\end{itemize}

\begin{table}[H]
\centering
\caption{FreeRTOS: Stärken und Risiken}
\label{tab:freertos}
\begin{tabularx}{\textwidth}{l X}
\toprule
\textbf{Stärke} & Klare Prioritäten, gute Trennung IO ↔ App \\
\textbf{Risiko} & Ownership/Synchronisation sauber halten (Regel~6) \\
\bottomrule
\end{tabularx}
\end{table}

% -----------------------------------------------------------------------------
\subsection{One-Hot-Ausgabe (Definition, byte-basiert)}
\label{sec:one-hot-ausgabe}

\textbf{Regel:} Erst löschen, dann genau 1~Bit setzen.

\begin{lstlisting}[style=arduino]
// Pseudo-Code
clear(led_bytes);
if (selected_id != 0) {
    uint8_t bit = selected_id - 1;
    led_bytes[bit / 8] |= (1u << (bit % 8));
}
\end{lstlisting}

\begin{infobox}[Passend zum verifizierten Mapping]
\begin{itemize}
  \item \code{OUT[0]}: LED1..LED8 → Bit0..Bit7
  \item \code{OUT[1]}: LED9..LED10 → Bit0..Bit1
\end{itemize}
\end{infobox}

% -----------------------------------------------------------------------------
\subsection{One-Hot-Checkliste (damit es wirklich immer stimmt)}
\label{sec:one-hot-checkliste}

\begin{itemize}
  \item[$\square$] \code{selected\_id} ist \textbf{ein} Wert (\code{0..10}), kein Bitfeld
  \item[$\square$] Active-low wird \textbf{einmal} normalisiert (\code{pressed = !raw\_bit})
  \item[$\square$] Event ist \textbf{Press-Edge} (released → pressed), nicht \glqq level\grqq
  \item[$\square$] LED-Ausgabe hat \textbf{einen zentralen Pfad}: \code{clear} → \code{set one bit}
  \item[$\square$] Reset/Boot setzt deterministischen Zustand (\zB \code{selected\_id=0}, alle LEDs aus)
  \item[$\square$] Debug/Logs: nur Events loggen (nicht jeden Scan)
\end{itemize}

% =============================================================================

\section{Desktop → Embedded Mindset (10 Regeln)}
\label{sec:embedded-mindset}

\begin{enumerate}
  \item \textbf{Hardware ist die API}\\
    Pins, Busse, Register sind \glqq die Schnittstelle\grqq. OS-Abstraktionen fehlen oft.

  \item \textbf{Worst-Case zählt, nicht der Durchschnitt}\\
    Plane Deadlines/Latenzen, nicht \glqq läuft meistens\grqq.

  \item \textbf{Kein Blockieren im kritischen Pfad}\\
    Keine langen ISR, keine Busy-Waits, kein \code{delay()} in Logik/Tasks.

  \item \textbf{ISR = minimal}\\
    ISR setzt Flag/Queue + Timestamp. Verarbeitung im Task/Mainloop.

  \item \textbf{Speicher ist ein Budget}\\
    Stack/RAM sind eng. Große Daten statisch oder in Pools. Rekursion vermeiden.

  \item \textbf{Nebenläufigkeit ist überall}\\
    ISR ↔ Task ↔ Main: Ownership und Synchronisation (Queue/Atomics) festlegen.

  \item \textbf{C++ Feature-Set bewusst wählen}\\
    \code{constexpr}, \code{enum class}, kleine POD-Structs: ja. Exceptions/RTTI/iostream: meist nein.

  \item \textbf{I/O entkoppeln}\\
    \code{read\_inputs()} → \code{logic\_update()} → \code{write\_outputs()} (testbar, deterministisch).

  \item \textbf{Fehler sind Normalfall}\\
    Timeouts, Brownout, Disconnect: Recovery + Watchdog einplanen.

  \item \textbf{Observability ist knapp}\\
    Logging sparsam, Counters, Debug-Pins/Tracing. Debugbarkeit mitdesignen.
\end{enumerate}

% -----------------------------------------------------------------------------
\subsection{Beispiel: Selection Panel (One-Hot) -- \glqq Taster i wählt LED i\grqq}
\label{sec:selection-panel-beispiel}

\textbf{Ziel:} 10~Taster, 10~LEDs. Bei Tastendruck wird \highlight{genau eine} LED aktiv (One-Hot). Ein neuer Tastendruck schaltet \textbf{sofort} auf die neue Auswahl um.

\subsubsection{Konventionen (wichtig, konsistent)}

\begin{itemize}
  \item \textbf{Buttons sind active-low (Hardware):} \code{raw\_bit = 0} bedeutet \glqq gedrückt\grqq{} (\zB CD4021).
  \item \textbf{Logik arbeitet normalisiert:} \code{pressed = !raw\_bit} → dann gilt in der Logik: \code{1 = gedrückt}.
  \item \textbf{Auswahl wird als ID geführt:}\\
    \code{selected\_id} $\in \{0..10\}$\\
    \code{0} = keine Auswahl, \code{1..10} = LED/Taster-ID
\end{itemize}

\subsubsection{Datenfluss (Architekturregel)}

\begin{lstlisting}[style=base,numbers=none]
read_inputs_raw() -> normalize_active_low() -> debounce_update()
                  -> selection_update() -> write_outputs_onehot()
\end{lstlisting}

\begin{table}[H]
\centering
\caption{Datenfluss-Funktionen im Detail}
\label{tab:datenfluss-funktionen}
\begin{tabularx}{\textwidth}{l X}
\toprule
\textbf{Funktion} & \textbf{Beschreibung} \\
\midrule
\code{read\_inputs\_raw()}      & Liest Bits aus dem Schieberegister (\zB 2~Bytes für 10~Taster) \\
\code{normalize\_active\_low()} & Invertiert Bits → Logik-Semantik \code{1 = pressed} \\
\code{debounce\_update()}       & Macht aus prellenden Rohwerten stabile Zustände \\
\code{selection\_update()}      & Setzt \code{selected\_id} bei neuem Press-Event \\
\code{write\_outputs\_onehot()} & Setzt genau \textbf{ein} LED-Bit (alle anderen 0) \\
\bottomrule
\end{tabularx}
\end{table}

% -----------------------------------------------------------------------------
\subsection{Arduino-Loop vs.\ FreeRTOS (Prinzip, nicht Code)}
\label{sec:loop-vs-rtos}

\subsubsection{Variante A: Arduino-Loop (ohne eigene Tasks)}

\code{loop()} läuft zyklisch (Polling), \zB alle \SI{1}{\milli\second}. Alle Schritte passieren sequentiell:

\begin{lstlisting}[style=base,numbers=none]
Inputs lesen -> normalisieren -> entprellen -> Event erkennen -> LED schreiben
\end{lstlisting}

\begin{table}[H]
\centering
\caption{Arduino-Loop: Stärken und Risiken}
\label{tab:arduino-loop}
\begin{tabularx}{\textwidth}{l X}
\toprule
\textbf{Stärke} & Simpel, wenige Nebenläufigkeitsprobleme \\
\textbf{Risiko} & \code{loop()} darf nicht blockieren (Regel~3) \\
\bottomrule
\end{tabularx}
\end{table}

\subsubsection{Variante B: FreeRTOS (Task-basiert)}

\begin{itemize}
  \item \textbf{IO-Task (high prio)} besitzt SPI und macht:\\
    Inputs lesen → normalisieren → entprellen → \code{selected\_id} bestimmen
  \item Optional: \textbf{App-Task} bekommt Events per Queue (\zB für Serial/Web)
\end{itemize}

\begin{table}[H]
\centering
\caption{FreeRTOS: Stärken und Risiken}
\label{tab:freertos}
\begin{tabularx}{\textwidth}{l X}
\toprule
\textbf{Stärke} & Klare Prioritäten, gute Trennung IO ↔ App \\
\textbf{Risiko} & Ownership/Synchronisation sauber halten (Regel~6) \\
\bottomrule
\end{tabularx}
\end{table}

% -----------------------------------------------------------------------------
\subsection{One-Hot-Ausgabe (Definition, byte-basiert)}
\label{sec:one-hot-ausgabe}

\textbf{Regel:} Erst löschen, dann genau 1~Bit setzen.

\begin{lstlisting}[style=arduino]
// Pseudo-Code
clear(led_bytes);
if (selected_id != 0) {
    uint8_t bit = selected_id - 1;
    led_bytes[bit / 8] |= (1u << (bit % 8));
}
\end{lstlisting}

\begin{infobox}[Passend zum verifizierten Mapping]
\begin{itemize}
  \item \code{OUT[0]}: LED1..LED8 → Bit0..Bit7
  \item \code{OUT[1]}: LED9..LED10 → Bit0..Bit1
\end{itemize}
\end{infobox}

% -----------------------------------------------------------------------------
\subsection{One-Hot-Checkliste (damit es wirklich immer stimmt)}
\label{sec:one-hot-checkliste}

\begin{itemize}
  \item[$\square$] \code{selected\_id} ist \textbf{ein} Wert (\code{0..10}), kein Bitfeld
  \item[$\square$] Active-low wird \textbf{einmal} normalisiert (\code{pressed = !raw\_bit})
  \item[$\square$] Event ist \textbf{Press-Edge} (released → pressed), nicht \glqq level\grqq
  \item[$\square$] LED-Ausgabe hat \textbf{einen zentralen Pfad}: \code{clear} → \code{set one bit}
  \item[$\square$] Reset/Boot setzt deterministischen Zustand (\zB \code{selected\_id=0}, alle LEDs aus)
  \item[$\square$] Debug/Logs: nur Events loggen (nicht jeden Scan)
\end{itemize}

% =============================================================================

\section{Desktop → Embedded Mindset (10 Regeln)}
\label{sec:embedded-mindset}

\begin{enumerate}
  \item \textbf{Hardware ist die API}\\
    Pins, Busse, Register sind \glqq die Schnittstelle\grqq. OS-Abstraktionen fehlen oft.

  \item \textbf{Worst-Case zählt, nicht der Durchschnitt}\\
    Plane Deadlines/Latenzen, nicht \glqq läuft meistens\grqq.

  \item \textbf{Kein Blockieren im kritischen Pfad}\\
    Keine langen ISR, keine Busy-Waits, kein \code{delay()} in Logik/Tasks.

  \item \textbf{ISR = minimal}\\
    ISR setzt Flag/Queue + Timestamp. Verarbeitung im Task/Mainloop.

  \item \textbf{Speicher ist ein Budget}\\
    Stack/RAM sind eng. Große Daten statisch oder in Pools. Rekursion vermeiden.

  \item \textbf{Nebenläufigkeit ist überall}\\
    ISR ↔ Task ↔ Main: Ownership und Synchronisation (Queue/Atomics) festlegen.

  \item \textbf{C++ Feature-Set bewusst wählen}\\
    \code{constexpr}, \code{enum class}, kleine POD-Structs: ja. Exceptions/RTTI/iostream: meist nein.

  \item \textbf{I/O entkoppeln}\\
    \code{read\_inputs()} → \code{logic\_update()} → \code{write\_outputs()} (testbar, deterministisch).

  \item \textbf{Fehler sind Normalfall}\\
    Timeouts, Brownout, Disconnect: Recovery + Watchdog einplanen.

  \item \textbf{Observability ist knapp}\\
    Logging sparsam, Counters, Debug-Pins/Tracing. Debugbarkeit mitdesignen.
\end{enumerate}

% -----------------------------------------------------------------------------
\subsection{Beispiel: Selection Panel (One-Hot) -- \glqq Taster i wählt LED i\grqq}
\label{sec:selection-panel-beispiel}

\textbf{Ziel:} 10~Taster, 10~LEDs. Bei Tastendruck wird \highlight{genau eine} LED aktiv (One-Hot). Ein neuer Tastendruck schaltet \textbf{sofort} auf die neue Auswahl um.

\subsubsection{Konventionen (wichtig, konsistent)}

\begin{itemize}
  \item \textbf{Buttons sind active-low (Hardware):} \code{raw\_bit = 0} bedeutet \glqq gedrückt\grqq{} (\zB CD4021).
  \item \textbf{Logik arbeitet normalisiert:} \code{pressed = !raw\_bit} → dann gilt in der Logik: \code{1 = gedrückt}.
  \item \textbf{Auswahl wird als ID geführt:}\\
    \code{selected\_id} $\in \{0..10\}$\\
    \code{0} = keine Auswahl, \code{1..10} = LED/Taster-ID
\end{itemize}

\subsubsection{Datenfluss (Architekturregel)}

\begin{lstlisting}[style=base,numbers=none]
read_inputs_raw() -> normalize_active_low() -> debounce_update()
                  -> selection_update() -> write_outputs_onehot()
\end{lstlisting}

\begin{table}[H]
\centering
\caption{Datenfluss-Funktionen im Detail}
\label{tab:datenfluss-funktionen}
\begin{tabularx}{\textwidth}{l X}
\toprule
\textbf{Funktion} & \textbf{Beschreibung} \\
\midrule
\code{read\_inputs\_raw()}      & Liest Bits aus dem Schieberegister (\zB 2~Bytes für 10~Taster) \\
\code{normalize\_active\_low()} & Invertiert Bits → Logik-Semantik \code{1 = pressed} \\
\code{debounce\_update()}       & Macht aus prellenden Rohwerten stabile Zustände \\
\code{selection\_update()}      & Setzt \code{selected\_id} bei neuem Press-Event \\
\code{write\_outputs\_onehot()} & Setzt genau \textbf{ein} LED-Bit (alle anderen 0) \\
\bottomrule
\end{tabularx}
\end{table}

% -----------------------------------------------------------------------------
\subsection{Arduino-Loop vs.\ FreeRTOS (Prinzip, nicht Code)}
\label{sec:loop-vs-rtos}

\subsubsection{Variante A: Arduino-Loop (ohne eigene Tasks)}

\code{loop()} läuft zyklisch (Polling), \zB alle \SI{1}{\milli\second}. Alle Schritte passieren sequentiell:

\begin{lstlisting}[style=base,numbers=none]
Inputs lesen -> normalisieren -> entprellen -> Event erkennen -> LED schreiben
\end{lstlisting}

\begin{table}[H]
\centering
\caption{Arduino-Loop: Stärken und Risiken}
\label{tab:arduino-loop}
\begin{tabularx}{\textwidth}{l X}
\toprule
\textbf{Stärke} & Simpel, wenige Nebenläufigkeitsprobleme \\
\textbf{Risiko} & \code{loop()} darf nicht blockieren (Regel~3) \\
\bottomrule
\end{tabularx}
\end{table}

\subsubsection{Variante B: FreeRTOS (Task-basiert)}

\begin{itemize}
  \item \textbf{IO-Task (high prio)} besitzt SPI und macht:\\
    Inputs lesen → normalisieren → entprellen → \code{selected\_id} bestimmen
  \item Optional: \textbf{App-Task} bekommt Events per Queue (\zB für Serial/Web)
\end{itemize}

\begin{table}[H]
\centering
\caption{FreeRTOS: Stärken und Risiken}
\label{tab:freertos}
\begin{tabularx}{\textwidth}{l X}
\toprule
\textbf{Stärke} & Klare Prioritäten, gute Trennung IO ↔ App \\
\textbf{Risiko} & Ownership/Synchronisation sauber halten (Regel~6) \\
\bottomrule
\end{tabularx}
\end{table}

% -----------------------------------------------------------------------------
\subsection{One-Hot-Ausgabe (Definition, byte-basiert)}
\label{sec:one-hot-ausgabe}

\textbf{Regel:} Erst löschen, dann genau 1~Bit setzen.

\begin{lstlisting}[style=arduino]
// Pseudo-Code
clear(led_bytes);
if (selected_id != 0) {
    uint8_t bit = selected_id - 1;
    led_bytes[bit / 8] |= (1u << (bit % 8));
}
\end{lstlisting}

\begin{infobox}[Passend zum verifizierten Mapping]
\begin{itemize}
  \item \code{OUT[0]}: LED1..LED8 → Bit0..Bit7
  \item \code{OUT[1]}: LED9..LED10 → Bit0..Bit1
\end{itemize}
\end{infobox}

% -----------------------------------------------------------------------------
\subsection{One-Hot-Checkliste (damit es wirklich immer stimmt)}
\label{sec:one-hot-checkliste}

\begin{itemize}
  \item[$\square$] \code{selected\_id} ist \textbf{ein} Wert (\code{0..10}), kein Bitfeld
  \item[$\square$] Active-low wird \textbf{einmal} normalisiert (\code{pressed = !raw\_bit})
  \item[$\square$] Event ist \textbf{Press-Edge} (released → pressed), nicht \glqq level\grqq
  \item[$\square$] LED-Ausgabe hat \textbf{einen zentralen Pfad}: \code{clear} → \code{set one bit}
  \item[$\square$] Reset/Boot setzt deterministischen Zustand (\zB \code{selected\_id=0}, alle LEDs aus)
  \item[$\square$] Debug/Logs: nur Events loggen (nicht jeden Scan)
\end{itemize}
